\documentclass[11pt]{article}
\usepackage[T1]{fontenc}
\usepackage{lmodern}
\usepackage{amsmath,amssymb,amsthm}
\usepackage{booktabs}
\usepackage[margin=1in]{geometry}
\usepackage[hidelinks]{hyperref}

\newtheorem{theorem}{Theorem}
\newtheorem{lemma}[theorem]{Lemma}
\newtheorem{proposition}[theorem]{Proposition}
\newtheorem{corollary}[theorem]{Corollary}
\theoremstyle{remark}
\newtheorem{remark}[theorem]{Remark}

\newcommand{\SigmaFour}{\{0,1,2,3\}}
\newcommand{\Xsf}{X_{\mathrm{sf},4}}
\newcommand{\Lang}{\mathcal{L}}
\newcommand{\Sym}{S_4}

\title{The Exact Locality Cost of Alphabet\\Symmetry in Square-Free Marker Rules}
\author{}
\date{Working manuscript --- 15 November 2026}

\begin{document}
\maketitle

\begin{abstract}
We determine optimal input-window radii for shift-equivariant Boolean
marker rules on the bi-infinite square-free subshift over four letters.
Selected sites must be at least three positions apart, and every
position must lie within distance two of a marker. Without imposed
alphabet symmetry the least symmetric coding radius is five; full
alphabet-permutation invariance under $S_4$ raises it to seven.
More precisely, for each subgroup $H\le S_4$, the optimum is five
if $|H|\le2$, six if $|H|>2$ and $H$ contains no four-cycle, and
seven if $H$ contains a four-cycle. These cases account for
10, 13 and 7 of the 30 subgroups, respectively.
Lower bounds combine finite packing/covering clauses realized
in genuine bi-infinite square-free sources, explicit extension
lemmas and independently checked propositional refutations.
Upper bounds use exhaustively validated finite-word marker
tables; an eleven-clause obstruction excludes four-cycle
invariance at radius six. We also classify all realizable exact
alphabet stabilizers at radius five. This fixed-parameter
classification quantifies the locality cost of alphabet
symmetry; it does not claim a new general marker-existence
theorem or a new SAT method.
\end{abstract}

\noindent\textbf{Keywords.}
square-free words; symbolic dynamics; local marker rules;
alphabet-permutation symmetry; coding radius; certified
finite-factor obstructions.

\noindent\textbf{2020 Mathematics Subject Classification.}
Primary 37B10; Secondary 68R15, 68R07.

\section{Introduction}

Classical marker lemmas provide finite-radius markers for aperiodic
zero-dimensional dynamical systems; see, for example, the modern
subshift formulation discussed by Meyerovitch
\cite{meyerovitch2025embedding}. Here the question is quantitative:
what is the least radius of a \emph{single} local rule for a prescribed
marker spacing and covering guarantee? A further constraint is
invariance under every renaming of the alphabet.

A concrete quantitative predecessor is T{\"o}rm{\"a}'s computable
$n$-marker block map, whose symmetric radius is bounded by
$n|A|^{2n+1}$ \cite[Lemma~6.10]{torma2015structural}.
At $n=3$ on a four-symbol alphabet this gives
$3\cdot4^7=49{,}152$. Its alternative to finding a marker within
distance $n-1=2$ is a central length-seven word with period less
than three. Such a word necessarily contains a square, so this
alternative is impossible on the square-free shift. Thus finite
radius and the desired packing/covering behavior are already
classically guaranteed; the sharp fixed-parameter radius is not.

Our focus is the quaternary square-free shift with
separation-three and domination-distance-two requirements.
Exact lower bounds need more than finite square-free
source words: the underlying constraints must be realized
in bi-infinite square-free configurations. Our proof
therefore separates finite SAT/DRAT evidence from the
Hall--Thue extension arguments, and checks positive rules
against every finite square-free source.

\paragraph{What is new in the fixed model.}
The minimum unrestricted radius is five.
Every single alphabet involution remains compatible with
that optimum, but larger subgroups require radius six,
and the presence of a four-cycle forces radius seven.
The full $S_4$ penalty is exactly two. A short
radius-five four-cycle collision and an eleven-clause
radius-six obstruction help explain the two levels.
The contribution is this quantitative subgroup
classification, not the classical existence of markers,
a new SAT method or a new graph-coloring algorithm.

\section{The model}

A finite word is \emph{square-free} if it contains no consecutive
repetition $uu$ with $u$ nonempty. Let
\[
  \Xsf=\{x\in\SigmaFour^{\mathbb Z}:
          \text{$x$ has no nonempty square factor}\}.
\]
Write $\Lang_n(\Xsf)$ for the set of length-$n$ factors occurring in
some $x\in\Xsf$.

A symmetric-radius-$r$ marker rule is a function
\[
 f:\Lang_{2r+1}(\Xsf)\longrightarrow\{0,1\}.
\]
At site $i$ on $x$, the indicator is
$m_f(x,i)=f(x_{i-r}\cdots x_{i+r})$. Validity means
\begin{align}
 m_f(x,i)=m_f(x,j)=1,\ i\ne j
   &\Longrightarrow |i-j|\ge3, \label{eq:separation}\\
 \forall i\in\mathbb Z\quad
   &\exists j\in[i-2,i+2]\colon m_f(x,j)=1.
   \label{eq:domination}
\end{align}
Equivalently, consecutive markers have gaps from three to five.
For instance, the set of sites
$\{8k,8k+3:k\in\mathbb Z\}$ has alternating gaps
three and five and obeys both output constraints.
This is only an illustrative marker pattern: it does
not define a single input-dependent block rule that
works on every square-free configuration.

The symmetric group $\Sym$ acts on $\Xsf$ by applying the same
alphabet permutation at every position. A rule is
$\Sym$-\emph{invariant} if
$f(\pi w)=f(w)$ for all $\pi\in\Sym$ and
$w\in\Lang_{2r+1}(\Xsf)$. Invariance under a subgroup $H\le\Sym$ is defined
by the same condition with $\pi\in H$. For such an
$H$, let $r_H(4,3)$ be the least symmetric coding
radius admitting a valid $H$-invariant rule; the
parameters $(4,3)$ denote respectively the alphabet
size and the required minimum separation between
markers (the covering distance is two throughout).
We write $r_{\mathrm{all}}(4,3)=r_{\{e\}}(4,3)$
for no imposed alphabet symmetry and
$r_{\Sym}(4,3)$ for the full group.

\begin{remark}
Write $\mathcal W_n$ for the finite set of all square-free words
of length $n$ over $\SigmaFour$. This set can strictly contain
$\Lang_n(\Xsf)$. For example,
$\texttt{010201030102010}\in\mathcal W_{15}$
is square-free but cannot be extended by
any one of the four symbols on either side,
so it does not belong to $\Lang_{15}(\Xsf)$.
Using $\mathcal W_n$ yields
legitimate sufficient conditions for an upper-bound construction, but
a lower-bound clause must be witnessed by a factor that really occurs
in a bi-infinite word. We maintain this distinction throughout.
\end{remark}

\subsection{Qualitative existence under alphabet symmetry}

Even before the quantitative construction below, a classical marker
lemma guarantees that full alphabet symmetry cannot obstruct marker
selection at \emph{every} radius. The following observation clarifies
which part of the result is genuinely quantitative.

\begin{proposition}[Qualitative symmetry-invariant markers]
\label{prop:existence}
For every finite group $G\le S_4$ acting by alphabet relabeling,
there exists a valid, $G$-invariant marker rule of some finite
symmetric coding radius on $\Xsf$.
\end{proposition}

\begin{proof}
The quotient $Y=\Xsf/G$ is compact and zero dimensional: since $G$
is finite, invariant clopen neighborhoods separate distinct
finite orbits. The shift descends to a homeomorphism of $Y$.
If $\sigma^p[x]=[x]$ for some $p\ge1$, then
$\sigma^p x=gx$ for some $g\in G$. Since $g$ has finite order
$h$, we obtain $\sigma^{ph}x=x$, contradicting the
square-freeness of $x$. Thus the quotient shift is aperiodic.

Apply a Krieger-type marker lemma
\cite[Lemma 3.4]{meyerovitch2025embedding} to $Y$
with the finite symmetric shift set $\{-2,-1,1,2\}$.
It yields a clopen set $C\subseteq Y$ with
$C\cap\sigma^d C=\varnothing$ for $d\in\{-2,-1,1,2\}$,
and whose translates for $d\in\{-2,-1,0,1,2\}$ cover
$Y$. Pulling $C$ back to $\Xsf$ gives
a clopen, $G$-invariant marker set. Since $\Xsf$ is a
compact finite-alphabet subshift, a clopen indicator is
determined by some finite symmetric coordinate window,
in accordance with the sliding-block principle
\cite{hedlund1969codes}.
Its shift translates define the desired shared rule.
\end{proof}

The argument provides no numerical coding-radius bound,
and it does not determine the minimum radius. In contrast,
the source-certified marker-table construction below gives
the exact full-$S_4$ radius seven.

\section{Main fixed-parameter statement}

\begin{theorem}[Exact locality cost of full alphabet symmetry]
\label{thm:main}
For the quaternary bi-infinite square-free shift
under \eqref{eq:separation}--\eqref{eq:domination},
\[
 \boxed{r_{\mathrm{all}}(4,3)=5,\qquad
        r_{\Sym}(4,3)=7.}
\]
Thus full alphabet-permutation invariance increases
the minimum symmetric coding radius by exactly two.
\end{theorem}
The unrestricted radius-four refutation and
the full-$S_4$ radius-six obstruction provide
the lower bounds. The radius-five unrestricted
and radius-seven fully invariant decision tables
furnish the matching upper bounds, verified below.

\section{A square-free two-sided extension lemma}

Fix one symbol, denoted $0$, as a separator. A finite
square-free source with a separator admits two different extension
arguments depending on whether that symbol occurs once or more
than once. Stating both cases is essential for the finite-source
lower-bound certificates.

\begin{lemma}[Single-separator extension]\label{lem:single-separator}
Let $Z$ be a finite square-free word over $\SigmaFour$
containing $0$ at exactly one position $a$. Suppose the entire
prefix $Z[0,a)$ is a terminal segment of a left-infinite
square-free ternary word over $\{1,2,3\}$, and the entire
suffix $Z(a,N)$ is an initial segment of a right-infinite
square-free ternary word over $\{1,2,3\}$, where
$Z(a,N)=z_{a+1}\cdots z_{N-1}$ and $N=|Z|$.
Then $Z$ occurs in a bi-infinite square-free word.
\end{lemma}
\begin{proof}
Identify the complete separator-free prefix and suffix
with their respective infinite ternary rays and glue them
to the single symbol $0$. A square containing $0$ would
have to contain the same symbol twice, once in each half,
which is impossible. A square avoiding $0$ lies wholly
within one of the two square-free ternary rays. Thus the
bi-infinite glued word contains no square.
\end{proof}

\begin{lemma}[Two-sided separator extension]\label{lem:separator}
Let $Z=z_0\ldots z_{N-1}$ be a finite square-free word over
$\SigmaFour$, with at least two occurrences of $0$. Let $a$ and $b$
be the indices of its first and last $0$, and write $D=b-a\ge1$.
Suppose
\[
   a\ge D-1,\qquad N-1-b\ge D-1.
\]
Suppose also that the \emph{entire} prefix $Z[0,a)$ is a terminal
segment of a left-infinite square-free word over $\{1,2,3\}$ and
the \emph{entire} suffix $Z(b,N)$ is an initial segment of a
right-infinite square-free word over $\{1,2,3\}$, where
$Z(b,N)=z_{b+1}\cdots z_{N-1}$.
Then $Z$ occurs in a bi-infinite square-free word over $\SigmaFour$.
\end{lemma}

\begin{proof}
Glue the chosen two ternary rays to $Z$, identifying the complete
zero-free prefix and suffix rather than adding duplicate copies.
A square lying entirely in $Z$ is impossible. A square containing no
$0$ that crosses a left or right outer join lies wholly within the
corresponding square-free ternary ray: it cannot cross the separator
block. Thus only a square containing $0$ and crossing an outer join
needs to be excluded.

Index $Z$ starting at $0$, with its left extension at negative
indices. Suppose a square of period $p$ starts at $s<0$ and contains
a $0$. Its right endpoint is $s+2p-1\le2p-2$. Let $t\ge a$ be its
first occurrence of $0$. This occurrence must lie in the first
half of the square, since otherwise the corresponding site at $t-p$
would be an earlier $0$. Therefore $t+p$ is another occurrence
of $0$ in the square, whence
\[
   p\le b-a=D.
\]
On the other hand $t+p$ lies before the right endpoint, so
$t+p\le 2p-2$ and
\[
   p\ge t+2\ge a+2\ge D+1,
\]
a contradiction. Reversal of indices gives the identical
contradiction for a square crossing the right join, using
$N-1-b\ge D-1$. Hence the glued word is square-free.
\end{proof}

The classical Hall--Thue word is obtained from the fixed point of
\[
  0\mapsto012,\qquad1\mapsto02,\qquad2\mapsto1.
\]
Its square-freeness goes back to Thue; the
specified Hall morphism and the historical proof are
described in \cite{badkobeh2022avoiding,lothaire2002words}.
The Hall morphism need not preserve square-freeness
on \emph{every} ternary square-free input: what is
used here is square-freeness of its \emph{fixed point}.
The substitution is primitive, so its shift orbit
closure is minimal and the fixed point is uniformly
recurrent \cite[Sec.~5.2]{queffelec2010substitution}.
Every finite factor therefore occurs in a
bi-infinite square-free orbit-closure point, with
arbitrarily long left and right extensions.

Some of the historical source checkers also allow
reversals of ternary Hall factors. This causes
no extra assumption: reversing an entire
bi-infinite square-free Hall orbit-closure point
preserves square-freeness and realizes every reversed
factor. Independently chosen alphabet relabelings
and (where used) reversals of the two zero-free
rays are permitted in
Lemmas~\ref{lem:single-separator}
and~\ref{lem:separator}; the proof only needs
each ray separately to be square-free.

\paragraph{Complete source-realization case distinction.}
The lower-bound certificate verifiers distinguish finite sources
by their extension construction. The written argument covers
every category:
\[
\begin{array}{@{}ll@{}}
\text{Finite witness class}&\text{Infinite extension argument}\\
\hline
\text{ternary Hall--Thue factor}&
  \text{two-sided Hall--Thue orbit closure}\\
\text{exactly one distinguished separator}&
  \text{Lemma~\ref{lem:single-separator}}\\
\text{at least two separators, with full ternary exteriors}&
  \text{Lemma~\ref{lem:separator}}
\end{array}
\]
The verifier classes named \texttt{single-separator} use
the second row (25, 30 and 31 original core sources
in the three-cycle, normal Klein-four and intransitive
Klein-four obstruction suites, respectively).
Classes called \texttt{multi-separator},
\texttt{padded-multi-separator},
\texttt{expanded-multi-separator} or
\texttt{external-double-separator}, and the corresponding
radius-four constructions, use the third row after
checking all separators and the complete ternary exteriors.
A source without a separator must be certified as a
Hall--Thue factor by the first row. The finite programs
check these premises; the infinite extension implications
are the mathematical lemmas above. No finite square-free
word is assumed automatically bi-infinitely extendible.

\section{Computer-assisted lower bounds}

Let $G\le S_4$ be either the trivial group (for unrestricted
rules) or an alphabet-permutation symmetry group. Write
$[w]_G$ for the orbit of a finite word $w$ under $G$ and
assign one Boolean variable $y_{[w]_G}$ to each length-$(2r+1)$
context orbit. For a finite word $v$, write $v[a,b]$ for
its inclusive coordinate subword with indices starting at zero.

The local conditions have the following exact propositional form.
For $d\in\{1,2\}$ and $v\in\mathcal W_{2r+1+d}$,
marker separation at displacement $d$ requires
\begin{equation}\label{eq:source-separation-clause}
  \neg y_{[v[0,2r]]_G}
  \ \lor\
  \neg y_{[v[d,d+2r]]_G}.
\end{equation}
For every $v\in\mathcal W_{2r+5}$, covering the center
position within distance two requires
\begin{equation}\label{eq:source-cover-clause}
  \bigvee_{d=0}^{4} y_{[v[d,d+2r]]_G}.
\end{equation}
Repeated orbit variables within a clause are interpreted in the
ordinary propositional sense.

\begin{proposition}[Finite-factor certificate semantics]
\label{prop:finite-clauses}
A satisfying assignment to all clauses
\eqref{eq:source-separation-clause}--\eqref{eq:source-cover-clause}
over the finite square-free word sets $\mathcal W_n$ defines
a valid, $G$-invariant radius-$r$ marker rule on $\Xsf$.
Conversely, every such valid rule satisfies the clauses
whose source words belong to the actual language
$\Lang_n(\Xsf)$. In particular, an unsatisfiable collection
of clauses each witnessed by a bi-infinitely realizable
source word excludes every valid radius-$r$ rule.
\end{proposition}
\begin{proof}
Assign to each source-language context the Boolean value of
its orbit variable. Every pair of sites at distance $d\in\{1,2\}$
is contained in a length-$(2r+1+d)$ factor; its separation
clause forbids selecting both sites. Every five consecutive
candidate marker sites have their radius-$r$ contexts inside
one length-$(2r+5)$ factor; its covering clause forces at
least one selection. This proves the first assertion, since
every factor of $\Xsf$ is finite and square-free. For the
converse, take any source word that occurs in a bi-infinite
configuration and apply the defining separation or covering
condition at the corresponding sites. A valid rule therefore
satisfies its associated clause. If these genuine-source
clauses are jointly inconsistent, no valid rule can exist.
\end{proof}

The two separation families have source lengths
$2r+2$ and $2r+3$; the covering family has source length
$2r+5$. Thus the unrestricted radius-four and radius-five
checks use lengths $(10,11,13)$ and $(12,13,15)$,
respectively, while the full-symmetry radius-six
obstruction uses $(14,15,17)$. This finite encoding
is standard; the theorem-critical additional issue for
an unsatisfiability proof is source realization in $\Xsf$.

\subsection{Unrestricted radius five}

The unrestricted result is proved by two complementary
finite certificates: a refuted radius-four subformula
and a fully checked radius-five assignment. At radius
four, there are 12,768 named square-free length-nine
context variables. The archived unsatisfiable core has
24,292 distinct clauses: 14,367 packing clauses
and 9,925 five-site covering clauses. This core
need not be a minimal UNSAT core; it is a fixed,
source-supported subformula, and checking its
unsatisfiability is enough.

The 11,981-clause first subformula has 101
direct Hall-factor certificates, 5,801
single-separator certificates, and 6,079
two-separator-with-margins certificates.
The remaining 12,311 clauses have 12,233
multiple-separator witnesses and 78
external-double-separator witnesses.
Two separately implemented auditors validate each
individual source word, its original variable/packing
or covering clause, its finite square-freeness and
the relevant complete Hall--Thue exteriors and
separator margins. Their union is checked to equal
the 24,292-clause archived core, with zero
unresolved source clauses.

For radius four, the ideal packing source lengths
are ten (distance one) and eleven (distance two).
The archived packing witnesses \emph{all} use
length-eleven words; a distance-one clause is
obtained from two adjacent length-nine subwindows
of that longer source. This does not change its
validity, provided the full length-eleven source
is bi-infinitely realizable, as the source auditors
verify. The original exact core and a separately
SHA-256-pinned DRAT trace are checked together
by the external DRAT proof checker. Thus the
mathematical source-realization lemmas and
the finite refutation jointly exclude radius four.

At radius five, a complete Boolean assignment
to all 87,936 named length-eleven finite
square-free contexts satisfies the packing
constraints on all 230,520 length-12 and
604,608 length-13 square-free words and
the covering constraints on all 4,156,392
length-15 square-free words. This is sufficient
even without knowing which of those finite
words extend bi-infinitely. Therefore
$r_{\mathrm{all}}(4,3)=5$, since a
smaller-radius rule can ignore additional
input coordinates.

\subsection{A ten-clause full-permutation obstruction at radius six}

For full $\Sym$-invariance at radius six, factor contexts related
by any alphabet relabeling must share one Boolean decision.
Our original DIMACS orbit-variable IDs are \emph{one-based}:
enumerate length-13 square-free words over $\SigmaFour$
lexicographically, retain those in canonical
first-appearance form (their distinct symbols first appear
as $0,1,2,3$), and assign consecutive IDs starting at one.
This gives 25,192 orbit variables. Every named length-13
context maps to its first-appearance normalization.
The length-14, -15 and -17 source families likewise enumerate
canonical square-free words lexicographically but use
\emph{zero-based} source ordinals; the original verifier
independently checks the precise mappings.
Exhaustively enumerated finite square-free source factors give a
source-identifiable constraint system. A reduced unsatisfiable core
contains ten original clauses involving eight Boolean orbit
variables. For readability, set
\[
\begin{array}{c|cccccccc}
\text{symbol}&A&B&C&D&E&F&G&H\\\hline
\text{orbit ID}&21696&4588&12753&7817&21695&12754&21697&21915
\end{array}
\]
and write the ten clauses as
\begin{align*}
&(\neg A\vee\neg B),\quad
 (\neg C\vee\neg D),\quad
 (\neg D\vee\neg E),\\
&(\neg F),\quad(\neg G),\quad(\neg C\vee\neg H),\\
&(C\vee B\vee F\vee G),\\
&(A\vee D\vee F\vee G),\\
&(E\vee C\vee F\vee G),\\
&(H\vee D\vee F\vee G).
\end{align*}

\begin{lemma}\label{lem:tenclauses}
These ten propositional clauses are jointly unsatisfiable.
\end{lemma}
\begin{proof}
The unit clauses force $F=G=0$. If $C=1$, then the clauses
$\neg C\vee\neg D$ and $\neg C\vee\neg H$ give $D=H=0$,
contradicting $H\vee D$. Thus $C=0$.
The clauses $E\vee C$ and $C\vee B$ now force $E=B=1$.
Then $\neg D\vee\neg E$ and $\neg A\vee\neg B$ force
$D=A=0$, contradicting $A\vee D$.
\end{proof}

The finite contradiction alone does not prove nonexistence on
$\Xsf$. Each of the ten source clauses must arise from an
actually occurring bi-infinite square-free configuration.
The archived source certificates fix the corresponding finite
source word and orbit identifiers. A separately audited extension
certificate gives a full ternary Hall--Thue factor for the
\emph{entire} zero-free extremity on each side, together with the
margins of Lemma~\ref{lem:separator}. The checker tests the complete
extremities, including symbols in the source word beyond its first
or last separator, rather than only checking appended padding.
Lemma~\ref{lem:separator} then establishes bi-infinite
extendibility. Together with Lemma~\ref{lem:tenclauses}, this
excludes every fully $\Sym$-invariant radius-six rule.
Monotonicity of the coding radius excludes smaller radii and yields
\[
 r_{\Sym}(4,3)\ge 7.
\]

Independent DIMACS/source checks, a direct check of all $2^8$
truth assignments, and the original and reduced DRAT proofs are
part of the archived evidence. The appendix identifies the ten source factors, while
\texttt{REPRODUCING.md} documents the exact verifier commands
and pinned proof artifacts.

\section{A supplementary full-symmetry coloring construction}

\begin{proposition}[Shortest proper alphabet-orbit color window]
\label{prop:minimum-color-window}
There is no $S_4$-invariant coloring of length-$L$
square-free factors, for $L\le13$, assigning different
colors to every pair of sites at distance two on
$\Xsf$. For $L=14$, a proper nine-color assignment
exists by the certified finite graph below.
Thus 14 is the smallest window length for this
\emph{proper-factor-coloring construction}, not a
lower bound on arbitrary marker rules.
\end{proposition}
\begin{proof}
The square-free length-15 word
\[
 V=\texttt{021320310213203}
\]
has first and third length-13 windows
$\texttt{0213203102132}$ and
$\texttt{1320310213203}$.
The latter is obtained from the former by the
alphabet four-cycle $0\mapsto1\mapsto2\mapsto3
\mapsto0$. They therefore must have the same
color under any $S_4$-invariant factor coloring,
contradicting distance-two properness.

This is a source factor of the bi-infinite shift:
the SHA-pinned orbit-preflight witness is a
square-free length-38 extension containing $V$
at offset 12. Its separator $0$ occurs at positions
12, 17, 20 and 25 (span $D=13$), with complete
ternary Hall--Thue exteriors of length $D-1=12$
on both sides. Lemma~\ref{lem:separator} gives
bi-infinite extension. Any hypothetical proper
coloring on a shorter word length could be applied
to a consistently aligned subword of each
length-13 window, yielding an impossible
length-13 proper coloring. The length-14 existence
claim is furnished by the separately checked
nine-color graph assignment below.
\end{proof}

Although not sharp, the radius-23 construction below
provides a distinct combinatorial marker method.
Its certificate does not depend on the new radius-seven
decision table.

For $w\in\mathcal W_{14}$, let $[w]$ be its orbit under
alphabet relabeling. A finite conflict graph has 66,047
length-14 square-free word orbits as vertices. Join two vertices
whenever their length-14 windows occur at displacement one or
two in a finite square-free word of length 15 or 16.
Exhaustive enumeration yields 173,183 length-15 and 453,998
length-16 canonical source rows and 627,159 distinct undirected
conflict edges. A stored nine-color assignment was independently
checked against every source conflict. This graph may include
finite words absent from $\Lang_n(\Xsf)$; enforcing the stronger
finite-word constraints suffices for an \emph{upper} bound without
assuming bi-infinite extendibility.

Let $c(i)\in\{0,\dots,8\}$ be the color of the orbit of the
centered factor $x[i-6,i+7]$. The independent graph audit ensures
\[
       c(i)\ne c(j)\qquad\text{when }1\le |i-j|\le2.
       \tag{*}\label{eq:proper}
\]
Define marker decisions by increasing color rank:
\[
 M_k(i)= [c(i)=k]\ \wedge\
       \bigwedge_{\substack{j<k\\d\in\{-2,-1,1,2\}}}
       \neg M_j(i+d),
       \quad 0\le k\le8,
\]
and select $i$ when $M(i)=\bigvee_{k=0}^8 M_k(i)$.
The recursion is over finitely many color ranks, not over
the spatial positions, and is therefore well defined on
bi-infinite configurations.

\begin{proposition}[Finite-color marker construction]
\label{prop:color}
If a shift-invariant color rule from centered length-$L$
windows uses $C$ colors and assigns different colors to
positions at distances one and two, the ordered-color
construction above yields marker separation at least three,
domination distance at most two, and symmetric radius at most
\[
  \left\lceil\frac{L-1}{2}\right\rceil+2(C-1).
\]
The marker rule inherits every alphabet symmetry of the
underlying color rule.
\end{proposition}

\begin{proof}
If selected sites at distance at most two have different
colors, the later-color site would be blocked by the earlier
selected site, a contradiction. Sites with equal colors
cannot lie that close by the hypothesis. Hence the separation
is at least three.

If a site of color $k$ was not selected, its defining
conjunction failed because a site of a smaller color had
already been selected within distance two. Every site is
therefore selected or dominated.

The decision $M_k(i)$ depends only on colors at sites
within $2k$ of $i$: each recursive call moves by at most two
sites and decreases the color rank by at least one. A
centered length-$L$ factor reads at most
$\lceil(L-1)/2\rceil$ positions in either direction.
Consequently the final decision has symmetric radius
at most $\lceil(L-1)/2\rceil+2(C-1)$. Color-orbit
invariance and the rank-only recursion preserve alphabet
symmetry; translation equivariance is immediate.
\end{proof}

For the actual length-14 color window $x[i-6,i+7]$,
the left and right symbol dependencies have radii six and
seven. The eight possible decreasing-rank steps add at most
16 positions on each side. Thus the final marker at $i$
depends only on $x[i-22,i+23]$, which is contained in
a symmetric radius-23 window. With $L=14$ and $C=9$,
Proposition~\ref{prop:color} gives
\[
            r_{\Sym}(4,3)\le7+2(8)=23.
\]
The independent coloring audit establishes properness;
the rank-dependency argument establishes locality.
This supplies a nonoptimal full-symmetry rule;
the sharp upper bound is checked separately.

\section{Alphabet-invariance subgroups at radius five}
\label{sec:subgroup}

At the exact unrestricted optimum, the allowable alphabet
symmetry can be classified. A subgroup $H\le S_4$ is
\emph{realizable at radius five} if some valid radius-five
rule is invariant under every $h\in H$ on
$\Lang_{11}(\Xsf)$. This describes which groups a rule
can respect, not necessarily its exact stabilizer.

\begin{lemma}[A short four-cycle obstruction]
\label{lem:fourcycle}
No valid radius-five marker rule on $\Xsf$ is invariant
under a four-cycle of the input alphabet.
\end{lemma}
\begin{proof}
Consider the square-free length-15 word
\[
 W=\texttt{012032130120321}.
\]
For the four-cycle $g$ given by
$0\mapsto2$, $1\mapsto0$, $2\mapsto3$, $3\mapsto1$,
the five successive length-11 windows have orbit
pattern $A,B,A,B,A$. Their first two representatives are
$\texttt{01203213012}$ and $\texttt{12032130120}$.
Indeed, applying the \emph{same} permutation $g$ takes
each window at position $i$ to that at $i+2$ for
$i=0,1,2$.
If the rule is invariant under $g$, separation at
distance two forces $A=B=0$, while domination at the
central site forces $A\vee B$, a contradiction.

The word $W$ really occurs in $\Xsf$: the explicitly
checked square-free extension
\[
 Z=\texttt{12101203213012032101202}
\]
contains it at offset three. The separator symbol $3$
occurs at positions 7, 10 and 15, giving span $D=8$;
the entire separator-free prefix and suffix both
have length $D-1=7$ and are factors of relabeled
ternary Hall--Thue rays. Renaming the separator and
applying Lemma~\ref{lem:separator} proves that
$Z$, hence $W$, is bi-infinitely extendible.
The bundled elementary checker directly verifies
these finite facts and all six conjugate four-cycles.
\end{proof}

\begin{theorem}[Radius-five subgroup classification]
\label{thm:subgroup}
A subgroup $H\le S_4$ is realizable at radius five if
and only if $|H|\le2$. In particular, every single
transposition and every single double transposition
is compatible with optimal locality, but no larger
alphabet-permutation subgroup is.
\end{theorem}
\begin{proof}
The separately archived positive SAT assignments,
independently checked against all named finite
square-free factors of lengths 11, 12, 13 and 15,
provide an $(0\,1)$-invariant rule and a different
$(0\,1)(2\,3)$-invariant rule. Alphabet conjugacy
gives existence for every order-two subgroup; the
trivial subgroup is also realizable.

For the converse, the certified three-cycle obstruction
and Lemma~\ref{lem:fourcycle} exclude every element
of order three or four from the invariance group $H$.
Thus each nonidentity element of $H$ has order two.
It follows that $H$ is abelian: since $a^2=b^2=(ab)^2=e$,
we have $ab=ba$ for any $a,b\in H$.

If $H$ contains a transposition, relabel it as $(0\,1)$.
Every element of $H$ then lies in its centralizer
\[
 \{e,(0\,1),(2\,3),(0\,1)(2\,3)\}.
\]
A group of order greater than two must be this entire
intransitive Klein-four group. If $H$ contains no
transposition, its nonidentity elements can only be
the three double transpositions; order greater than
two then forces the normal Klein-four group.
The independent finite-source, two-sided Hall--Thue
extension and DRAT certificates exclude \emph{both}
Klein-four types at radius five. Hence $|H|\le2$.
\end{proof}

The negative three-cycle, normal Klein-four, and
intransitive Klein-four certificates respectively
verify 12,747, 53,487, and 26,819 source-supported
clauses with standalone UNSAT proofs. The separate
positive witnesses are checked over all original
named finite contexts, not only orbit representatives.
The finite-to-bi-infinite step for the negative cases
uses the classical Hall--Thue recurrence facts and
Lemma~\ref{lem:separator}; exact clean-checkout
verification commands appear in \texttt{REPRODUCING.md}.

Theorem~\ref{thm:subgroup} classifies which subgroups
can already be imposed at radius five. The next
theorem determines the exact minimum radius for
every subgroup. The radius-five exact-stabilizer
classification remains Corollary~\ref{cor:exact-stabilizers}.

\section{Exact locality for every alphabet subgroup}
\label{sec:exact-radii}

\begin{lemma}[Eleven-clause radius-six four-cycle obstruction]
\label{lem:new-c4-r6}
No valid radius-six rule on $\Xsf$ is invariant
under an alphabet four-cycle.
\end{lemma}
\begin{proof}
Take $g=(0\,1\,2\,3)$ and let $A,\dots,I$ name
nine distinct orbits of length-13 factors under
$\langle g\rangle$. The eleven source words of lengths
14, 15 and 17, and their exact packing and covering
clauses, appear in Table~\ref{tab:c4-core}.
Each source word has a checked finite square-free
extension with at least two distinguished separators.
The complete ternary exteriors are relabeled
\emph{forward} Hall--Thue factors; their lengths
meet the separator span-minus-one margin condition.
Lemma~\ref{lem:separator} therefore makes every
source bi-infinitely realizable.

The two unit clauses force $D=F=0$.
If $C=1$, the clauses
$\neg C\vee\neg G$ and $\neg C\vee\neg H$
would force $G=H=0$, contradicting
$D\vee F\vee G\vee H$. Hence $C=0$.
The clauses $B\vee C\vee D\vee F$
and $C\vee D\vee E\vee F$ now force
$B=E=1$, so $G=I=0$. The two remaining
covering clauses force $H=A=1$, contrary to
$\neg A\vee\neg H$.

The independent small-core auditor reconstructs
the orbits and clauses from all eleven source
words, verifies their complete Hall--Thue
extremities and directly rejects all $2^9$
Boolean assignments. Every four-cycle is a
conjugate of $g$, so the obstruction holds
for each such element.
\end{proof}

\begin{theorem}[Exact subgroup-radius classification]
\label{thm:exact-subgroup-radii}
For every alphabet-permutation subgroup $H\le S_4$,
the smallest symmetric radius of a valid
$H$-invariant marker rule is
\[
\boxed{
r_H(4,3)=
\begin{cases}
5,& |H|\le2,\\
6,& |H|>2\text{ and $H$ has no four-cycle},\\
7,& \text{$H$ contains a four-cycle}.
\end{cases}}
\]
Thus, among the 30 subgroups of $S_4$,
respectively 10, 13 and 7 have minimum radii
five, six and seven.
\end{theorem}
\begin{proof}
Theorem~\ref{thm:subgroup} establishes the
first case and excludes radius five for larger
groups. Lemma~\ref{lem:new-c4-r6} excludes
radius six for a group with a four-cycle.
Any still smaller rule could ignore unused
input letters and hence would also give a
radius-six rule, which is impossible.

For the upper bounds, four SHA-pinned finite
marker-decision tables give a radius-six
$A_4$-invariant rule, a radius-six
point-stabilizer $S_3$-invariant rule,
a radius-six intransitive $V_4$-invariant
rule, and a radius-seven $S_4$-invariant
rule. A C++17 checker exhaustively validates packing at
distances one and two and five-site covering on
finite square-free sources modulo the respective
alphabet group. A separate C11 implementation
expands each table to all named square-free
source words and checks the same constraints
without orbit pruning. These stronger finite-word
checks suffice for each upper bound without assuming
finite square-free words extend bi-infinitely.

Every subgroup of order greater than two
without a four-cycle is, up to conjugacy,
contained in $A_4$, a point stabilizer $S_3$,
or an intransitive $V_4$. An independent
enumerator checks this group reduction
against all 30 subgroups of $S_4$.
If $f$ is a valid $H$-invariant rule,
then $f^\pi(w)=f(\pi^{-1}w)$ is a valid
$\pi H\pi^{-1}$-invariant rule.
The fully invariant radius-seven rule works
for every remaining subgroup. All three
cases follow.
\end{proof}

\begin{corollary}[Exact radius-five alphabet stabilizers]
\label{cor:exact-stabilizers}
For valid radius-five marker rules, the realizable
\emph{exact} alphabet stabilizers belong to precisely
three conjugacy types:
\[
 \{e\},\qquad \langle(0\,1)\rangle,\qquad
 \langle(0\,1)(2\,3)\rangle.
\]
In particular, some valid radius-five rule has trivial
exact alphabet stabilizer.
\end{corollary}
\begin{proof}
Theorem~\ref{thm:subgroup} rules out stabilizers of
order greater than two. The two positive single-involution
SAT witnesses already realize both order-two conjugacy
types; they cannot have larger stabilizers.

For the trivial type, take the original source-certified
unrestricted radius-five SAT assignment. Its first audit
reconstructs all 87,936 finite square-free length-11
contexts and checks all 24 alphabet permutations. Only
the identity preserves that \emph{finite} decision table;
the nearest nonidentity permutation changes 10,528
decisions. This alone is not a proof about $\Xsf$, since
some finite square-free contexts may fail to be
bi-infinitely extendible.

A separate verifier closes this logical gap directly.
It computes the 36 distinct length-11 ternary factors
in a length-4096 prefix of the classical Hall--Thue
fixed point and all their four-letter alphabet relabelings,
giving 432 distinct quaternary contexts. Each is in
$\Lang_{11}(\Xsf)$: by the square-freeness and
uniform recurrence of the primitive Hall--Thue
substitution, every such ternary factor lies in a
bi-infinite square-free orbit-closure point; applying
any alphabet relabeling to that entire point preserves
square-freeness. Among these 432 genuine-language
contexts, every one of the 23 nonidentity alphabet
permutations changes at least one decision of the
original model. The smallest nonidentity disagreement
count on this set is 48, attained by the relabeling
with image word $\texttt{0213}$ (which swaps $1$ and $2$).
For example, this permutation sends
$\texttt{01020121020}$ to $\texttt{02010212010}$,
where the archived model selects the first context
and rejects the second. Both are relabeled factors
of the Hall--Thue word. The verifier archives an
explicit Hall factor, occurrence offset and
relabeling for a conflicting context for \emph{each}
nonidentity permutation.

Thus no nonidentity permutation fixes the model
\emph{on the actual shift}. Its separately certified
marker validity then proves the stated trivial exact
stabilizer.
\end{proof}

\paragraph{Results at a glance.}
Table~\ref{tab:main-results} distinguishes
minimum radii from exact stabilizer types.
\begin{table}[htbp]
\centering
\small
\begin{tabular}{@{}p{2.02in}p{0.65in}p{2.80in}@{}}
\toprule
Alphabet invariance & Radius & Evidence\\
\midrule
$|H|\le2$ & $5$ &
Radius-four obstruction; unrestricted and
two involution-type radius-five models\\
$|H|>2$, no four-cycle & $6$ &
Radius-five subgroup obstructions;
$A_4$, $S_3$, intransitive $V_4$
radius-six decision tables\\
$H$ contains a four-cycle & $7$ &
Eleven source-supported clauses at radius six;
fully $S_4$-invariant radius-seven table\\
\midrule
Exact stabilizers at radius five & $\{e\}$ or $C_2$ &
Three certified models, including
Hall-backed actual-language witnesses\\
\bottomrule
\end{tabular}
\caption{Complete fixed-case classification
(alphabet size four, separation three,
domination radius two). The last row is
not an additional radius.}
\label{tab:main-results}
\end{table}

\paragraph{Proof dependency roadmap.}
The computer-assisted \emph{lower} bounds have three distinct
links: (i) an independently reconstructed finite source clause
corresponds to a required packing or covering condition;
(ii) the source word is realized in the bi-infinite square-free
shift using its complete Hall--Thue exteriors and the appropriate
single- or multiple-separator lemma; and (iii) an independently
checked refutation shows the supported clauses are incompatible.
The \emph{upper} bounds instead check complete
marker-decision tables or an alternative orbit
factor coloring against \emph{all} finite square-free
sources. The nonoptimal coloring additionally
uses the rank-dependency argument.
No finite sample test substitutes for the infinite extension
arguments.

\section{Relationship to prior work and limitations}

\paragraph{Topological and explicit marker existence.}
Krieger-type marker lemmas give clopen sets with disjoint selected
translates and covering properties for aperiodic zero-dimensional
dynamical systems. Meyerovitch's
\cite[Lemmas~3.3--3.4]{meyerovitch2025embedding}
formulation covers subshifts and general compact totally
disconnected actions. Applied to the finite alphabet-permutation
quotient of our square-free shift, it already guarantees
\emph{qualitative} $S_4$-invariant markers at some finite
symmetric coding radius.

Importantly, \emph{the separation and covering requirements
themselves are classical}. On the free part of the binary
$\mathbb Z$-Bernoulli shift, Gao and Jackson
\cite[Lemma~2.1]{gaojackson2015markers}
construct a relatively clopen marker set whose points are
separated by more than $d$ and whose orbits are covered
within distance $d$. Taking $d=2$ gives precisely
separation at least three and covering within two;
their Lemma~2.9 gives stronger gap control.
This is a theorem on a noncompact free part: relative
clopenness there does not itself assert any \emph{uniform}
finite coding radius for a block rule on our compact
quaternary square-free shift, nor input-alphabet
$S_4$-invariance.

Macdonald's restatement of Krieger's cylinder-marker
lemma \cite[Lemma~2.6]{macdonald2024encoding}
gives markers separated by $N$, with a short-periodicity
alternative in a missing-marker neighborhood.
That alternative is excluded on a square-free subshift
for $N=3$. T{\"o}rm{\"a}'s explicit marker map
\cite[Lemma~6.10]{torma2015structural}
has radius at most $n|A|^{2n+1}$, specializing to
$49{,}152$ when $n=3$ and $|A|=4$.
It supplies neither optimality nor full alphabet
relabeling invariance.

Gao and Wang's stronger clopen marker existence theorems
\cite[Theorems~1.2--1.3]{gao2026markers}
and their later arbitrary-generating-set extension
\cite[Theorem~1.2]{gaowang2026arbitrary}
also work on free Bernoulli actions and allow a
covering constant chosen with the separation parameter.
They do not state the minimum fixed uniform input-window
radius with domination distance two and $S_4$-invariance
on $\Xsf$. Our contribution is the \emph{exact coding-radius
and alphabet-subgroup classification in this fixed system},
not the existence of separated and covering clopen
markers or the generic marker-construction mechanism.

\paragraph{Local algorithms and ruling sets.}
Our marker outputs form a $(3,2)$-ruling set on the line:
chosen positions have mutual distance at least three,
and every position is within distance two of one.
This is established terminology; for example,
Bourreau et al.\ give exactly these two properties in
\cite[Definition~6]{bourreau2026ruling}. Their
labeled-line rendezvous algorithm, however, depends
on node labels and distance, rather than giving the
least uniform radius of one alphabet-$H$-invariant
Boolean block predicate on our square-free shift.
Naor and Stockmeyer's LOCAL/LCL framework
\cite{naor1995local}, path-LCL decidability
\cite{balliu2019paths}, and automata-theoretic
classification \cite{chang2023automata} address
broader distributed computational models.
Even \emph{order invariance} of processor identifiers
is not invariance under all permutations of four
input symbols. Classical deterministic color
reduction \cite{colevishkin1986} and continuous
combinatorial colorings \cite{bernshteyn2023probabilistic}
are likewise background, not algorithms claimed as new.

\paragraph{Local string sampling and synchronizing sets.}
String synchronizing sets
\cite{kempa2019synchronizing,kociumaka2024ipm}
are a closer selection analogy. The precise
definition in \cite[Definition~5.1]{ellert2026synchronizing}
requires matching length-$2\tau$ contexts to
have consistent membership, and a marker in every
$\tau$-site window except in sufficiently periodic
regions. At $\tau=3$, square-freeness rules out the
period-one exception, so the coverage resembles
our problem. Crucially, it imposes \emph{no
minimum distance between selected sites}; aggregate
size bounds are not a hard packing guarantee.
Moreover, per-text consistency does not by itself
supply a single shared block table valid uniformly
on every point of $\Xsf$.

Minimizers and winnowing
\cite{schleimer2003winnowing,shaw2022localselection}
also guarantee a representative in every
designated window, but may select nearby positions.
Universal hitting sets \cite{zheng2021uhs}
give an analogous unavoidable-factor coverage
condition without our separation requirement.
Common minimizers use an ordering of $k$-mers;
the separate alphabet-reordering optimization
studied in \cite{verbeek2026reordering} does not
make their output invariant under every renaming.
None of these properties alone entails the exact
radius-five/six/seven $S_4$-subgroup classification.
We do not assert an improvement in string-indexing
running time, sampling density, or sequence matching.

\paragraph{Alphabet-permutation symmetry and color blindness.}
Salo and T{\"o}rm{\"a} introduced a systematic
theory of \emph{color-blind} cellular automata,
including invariance under specified symbol
permutations, in \cite{salotorma2014colorblind}.
Their Lemma~2 characterizes permutation-invariant
sets of finite words by equality/inequality
patterns, and their Proposition~1 addresses
permutation-equivariant cellular automata on
subshifts. Thus neither permutation-invariant
local descriptions nor symbol-orbit reduction
is introduced here. Their principal maps return
symbols of the \emph{input} alphabet, whereas
our rule returns a Boolean marker indicator
on a designated square-free subshift; most
importantly, their results do not state our
minimum \emph{uniform input radius} under
packing, covering, and each $H\le S_4$.

In particular, \emph{shift equivariance} means
all sites use one translation-invariant predicate,
while \emph{alphabet-$S_4$ invariance}
additionally forbids exploiting input symbol names.
This differs from having unique processor
identifiers in distributed LOCAL algorithms.
We do not claim the generic construction of
maximal independent sets by processing colors
in order, or the use of orbit representatives,
as new. Our contribution is the exact
fixed-parameter subgroup-radius classification
with source-backed matching lower bounds.

\paragraph{Square-free words and certificate extendibility.}
The existence and square-freeness of the ternary
Hall--Thue word are classical; the particular
substitution is documented in
\cite{badkobeh2022avoiding}. Square-free
constructions have also played a role in
distributed-algorithm lower bounds since
Attiya, Snir and Warmuth's anonymous-ring
work \cite{attiya1988anonymous}. That paper
does not solve our fixed-radius local marker
problem, but this historical connection
deserves acknowledgment.
Grytczuk, Kordulewski and Pawlik study square-free word
extension phenomena and extremal words
\cite{grytczuk2021squarefree}; thus finite square-freeness
by itself must not be mistaken for infinite extendibility.
Our proof does not infer the existence of a bi-infinite
continuation from finite square tests alone. It uses complete zero-free extremities and
Lemma~\ref{lem:separator}, followed by classical recurrence.

Our candidate contribution is the exact
two-unit symmetry locality penalty,
all 30 subgroup-minimum radii, and
the exact radius-five stabilizer types.
Our targeted primary-source comparisons
have not identified a theorem giving
these fixed-case values. We do not claim
universal bibliographic priority or an
all-parameter theorem, nor any
data-chunking performance improvement.

\section{Reproducibility and certificate scope}

The archived computations consist of an unrestricted
radius-four refutation, an unrestricted radius-five
construction, a ten-source full-$\Sym$ radius-six refutation
with corrected two-sided Hall--Thue source extensions, and
an eleven-source cyclic four-group radius-six
obstruction, four additional positive subgroup
marker tables, and an independently checked
nonoptimal nine-color factor graph. The
radius-four and subgroup finite-source audits also include
single-separator witnesses, covered by
Lemma~\ref{lem:single-separator}. Original
certificate byte lengths and SHA-256 values are recorded in
the artifact manifest.

The 21 original main-theorem source artifacts have their
byte counts and SHA-256 digests checked against a fixed manifest.
The subgroup annex separately checks the imported original Git blobs,
regenerates and audits the two involution quotient CNFs, verifies the
archived positive SAT models on all named finite factors, and
independently replays the source-certified C3 and both Klein-four
refutations. A source-pinned build of the MIT-licensed third-party
DRAT-trim tool checks the archived refutations.
For the unrestricted radius-four core, the full-$S_4$
ten-clause pilot, and three subgroup negative cores,
converted LRAT proofs were also accepted by the
separately formally verified CakeML \texttt{cake\_lpr}
checker. This independently confirms propositional
UNSAT, not the mathematical source-realization
premises. The original nine-job proof suite and the
separate exact-symmetry annex have passed in clean
private GitHub Actions checkouts. The final source
release still requires licensing and public-content
review before a fresh-history public artifact can
be prepared.

The checker and the Python finite-factor programs provide
independently executable evidence, but the classical infinite-word
facts and the two-sided extension argument remain mathematical
ingredients of the proof, not conclusions of finite execution.
The targeted primary-source originality comparison
and 26-entry bibliography reconciliation are recorded
in the private review ledger; they do not establish
exhaustive publication priority. This draft has not
undergone external peer review.

\paragraph{Proof and code availability (private review version).}
The 21 original main-theorem source artifacts are identified by
byte lengths and SHA-256 digests in
\texttt{artifacts/manifest.json}, while four positive
symmetry tables and the eleven-source four-cycle core
are identified in
\path{artifacts/exact-symmetry-addendum-manifest.json}.
The certificate versions are unchanged from private proof
baseline \texttt{59a7075f75b4}; exact reproduction commands,
checker source revisions and trust boundaries are documented
in \texttt{REPRODUCING.md}. The source and proof archive
is currently private, with no assigned public archive DOI
or reproducibility download URL. Public distribution requires
a separate provenance, third-party rights and licensing review.
A public source revision and persistent location will be
specified only when such a release is authorized.

\begingroup
\small
\bibliographystyle{plain}
\bibliography{references}
\endgroup

\clearpage
\appendix
\section{The ten finite source factors}

Table~\ref{tab:sources} lists the source words in the order of the
ten clauses above. Here $D$ is the distance from the first to
the last distinguished $0$ in the factor. For each row, the
separately checked left and right strings in
Table~\ref{tab:pads} extend the source to a square-free
finite witness. The entire zero-free extremities of that
witness (not necessarily these added strings alone) are
independently checked to be factors of relabeled forward
Hall--Thue words.

\begin{table}[ht]
\centering
\small
\begin{tabular}{rcl}
\toprule
Row & $D$ & Canonical square-free source factor\\
\midrule
1 & 10 & \texttt{01023013210231}\\
2 & 13 & \texttt{01201302312010}\\
3 & 13 & \texttt{012013023120103}\\
4 & 11 & \texttt{012032130120321}\\
5 & 13 & \texttt{012310320123103}\\
6 & 9 & \texttt{012312013023121}\\
7 & 13 & \texttt{01023013210230131}\\
8 & 14 & \texttt{01201302312013021}\\
9 & 15 & \texttt{01231032012310302}\\
10 & 14 & \texttt{01231201302312013}\\
\bottomrule
\end{tabular}
\caption{Original source factors of the ten unsatisfiable clauses.}
\label{tab:sources}
\end{table}

\begin{table}[ht]
\centering
\scriptsize
\begin{tabular}{rll}
\toprule
Row & Left appended ternary word & Right appended ternary word\\
\midrule
1 & \texttt{121323123} & \texttt{321232131}\\
2 & \texttt{121312321323} & \texttt{212312132123}\\
3 & \texttt{121312321323} & \texttt{121312321312}\\
4 & \texttt{1213231232} & \texttt{2321312313}\\
5 & \texttt{121312321323} & \texttt{121312321312}\\
6 & \texttt{12131232} & \texttt{32123121}\\
7 & \texttt{121312321312} & \texttt{231321312321}\\
8 & \texttt{1213123213121} & \texttt{2321312313213}\\
9 & \texttt{12131232131213} & \texttt{12312132123132}\\
10 & \texttt{1213123213121} & \texttt{2123213123132}\\
\bottomrule
\end{tabular}
\caption{Finite extensions verified against complete Hall--Thue tails.}
\label{tab:pads}
\end{table}

The ten original CNF source-clause identities and both full-extremity
audit reports are machine-checkable in the proof artifact package.
The graph-coloring certificate has its own independent source-edge
auditor. These are finite computations: the passage to infinite
configurations uses Lemma~\ref{lem:separator} and the cited classical
Hall--Thue theorem, not an extrapolation from finite test length.

\clearpage
\section{The eleven-source radius-six four-cycle core}
For $g=(0\,1\,2\,3)$, the nine
$\langle g\rangle$-context orbits are denoted
$A,\ldots,I$. The independent verifier recovers
the orbit-to-letter association directly from
the source windows; labels alone are not proof.
Each source has a complete SHA-pinned square-free
extension with forward Hall--Thue exteriors.

\begin{table}[htbp]
\centering
\small
\begin{tabular}{@{}rll@{}}
\toprule
Row & Square-free source word & Clause\\
\midrule
1 & \texttt{01031021320310213} & $A\vee D\vee F\vee I$\\
2 & \texttt{010310213203132} & $\neg A\vee\neg H$\\
3 & \texttt{01203102132031021} & $B\vee C\vee D\vee F$\\
4 & \texttt{012031021320313} & $\neg B\vee\neg G$\\
5 & \texttt{01320310213202} & $\neg C\vee\neg G$\\
6 & \texttt{013203102132021} & $\neg C\vee\neg H$\\
7 & \texttt{01320310213203101} & $C\vee D\vee E\vee F$\\
8 & \texttt{021320310213203} & $\neg D$\\
9 & \texttt{031021320310213} & $\neg F$\\
10 & \texttt{03102132031021310} & $D\vee F\vee G\vee H$\\
11 & \texttt{03203102132030} & $\neg E\vee\neg I$\\
\bottomrule
\end{tabular}
\caption{An eleven-clause radius-six obstruction
for a four-cycle of alphabet names.}
\label{tab:c4-core}
\end{table}

\end{document}
